\documentclass[aspectratio=169,11pt,svgnames]{beamer}

\usepackage[czech]{babel}
\usepackage[czech=quotes]{csquotes}
\usepackage{graphicx}
\usepackage{enumitem}
\usepackage{amsmath}
\usepackage{mathtools}
\usepackage{float}
\usepackage{tikz}

% Flowchart stuff
\usetikzlibrary{shapes.geometric, arrows.meta, calc, positioning}
\tikzstyle{startstop} = [rectangle, rounded corners, minimum width=3cm, minimum
height=1cm,text centered, draw=black, fill=Red!30]
\tikzstyle{io} = [trapezium, trapezium left angle=70, trapezium right angle=110,
minimum width=3cm, minimum height=1cm, text centered, draw=black, fill=Blue!30]
\tikzstyle{process} = [rectangle, minimum width=3cm, minimum height=1cm, text
centered, draw=black, fill=Orange!30]
\tikzstyle{decision} = [diamond, aspect=2, minimum width=3cm, minimum height=.5cm, text
centered, draw=black, fill=Green!30]
\tikzstyle{connector} = [draw, -latex']

\usepackage{pgfopts}
\usepackage{xcolor}
\usepackage{tcolorbox}
\usepackage{booktabs}

\usepackage[czech]{algorithm2e}

\usetheme[
 titlestyle=style2,
 titleformat=smallcaps,
 sectionstyle=plain,
 slidestyle=cyber,
 headingcolor=theme,
 block=transparent
]{trigon}

\title{Programovací jazyky}
\date{\today}
\author{Adam Klepáč}
\institute[GEVO]{Gymnázium Evolution Jižní Město}
\biglogo[width=.2\textwidth]{logo}
\smalllogo[width=.1\textwidth]{logo}
\titlegraphic{\includegraphics[height=\paperheight]{title}}

\def\subsectionname{}

% enumerate global settings
\setlist[enumerate,1]{label=\arabic*.}
\setlist[enumerate,2]{label=\alph*)}
\setlist[itemize,1]{label=\textbullet}
\setlist[itemize,2]{label=\textopenbullet}

% custom colors %
\definecolor{Night}{HTML}{100D14}
\definecolor{RobinEggBlue}{HTML}{05C5C8}
\definecolor{OxfordBlue}{HTML}{0C152F}
\definecolor{Plum}{HTML}{8B4184}
\definecolor{CoralPink}{HTML}{E18175}

\colorlet{tPrim}{RobinEggBlue}
\colorlet{tSec}{OxfordBlue}
\colorlet{tAccent}{Plum}
\colorlet{tTheme}{RobinEggBlue}
\colorlet{tTxt}{Night}

\tcbset{
 boxsep=7pt,
 fonttitle=\sc,
 colframe=tGreyBg,
 colframe=tTheme,
 boxrule=1pt
}

\begin{document}
\titleframe

\begin{frame}
 \frametitle{Obsah}
 \tableofcontents
\end{frame}

\section{Počítačový program}

\begin{frame}
 \frametitle{Co je počítačový program?}
 Počítačový program je sada instrukcí (uložená ve vnitřní paměti), kterou má CPU
 vykonat.\\
 \pause
 Popis a způsob zápisu instrukcí, které CPU umí vykonat, je určen jeho
 \emph{instrukční sadou}.\\
 \pause
 Instrukce společné většině instrukčních sad jsou:
 \begin{itemize}
  \item \alert{Datové a paměťové}
  \begin{itemize}
   \item \alert{SET} -- nastaví registr na konkrétní hodnotu;
   \pause
   \item \alert{COPY} (též \alert{MOVE}) -- zkopíruje data z registru nebo místa
    v paměti do registru nebo místa v paměti;
   \pause
   \item \alert{READ}/\alert{WRITE} -- přečti/zapiš data z připojeného/na
    připojené zařízení.
  \end{itemize}
 \end{itemize}
\end{frame}

\begin{frame}
 \frametitle{Co je počítačový program?}
 Počítačový program je sada instrukcí (uložená ve vnitřní paměti), kterou má CPU
 vykonat.\\
 Popis a způsob zápisu instrukcí, které CPU umí vykonat, je určen jeho
 \emph{instrukční sadou}.\\
 Instrukce společné většině instrukčních sad jsou:
 \begin{itemize}
  \item \alert{Výpočetní a logické}
  \begin{itemize}
   \item \alert{ADD}/\alert{SUBTRACT}/\alert{MULTIPLY}/\alert{DIVIDE} -- provede
    danou aritmetickou operaci na hodnoty ve dvou registrech a výsledek uloží
    opět do registru;
   \pause
   \item \alert{AND}/\alert{OR}/\alert{NOT} -- provede danou logickou operaci na
    hodnoty ve dvou registrech (v případě \alert{AND} a \alert{OR}) nebo jednom
    registru (v případě \alert{NOT}) a výsledek uloží do registru;
   \pause
   \item \alert{COMPARE} -- porovná hodnoty ve dvou registrech a výsledek uloží
    do registru;
   \pause
   \item instrukce související s aritmetikou plovoucího bodu (floating-point).
  \end{itemize}
 \end{itemize}
\end{frame}

\begin{frame}
 \frametitle{Co je počítačový program?}
 Počítačový program je sada instrukcí (uložená ve vnitřní paměti), kterou má CPU
 vykonat.\\
 Popis a způsob zápisu instrukcí, které CPU umí vykonat, je určen jeho
 \emph{instrukční sadou}.\\
 Instrukce společné většině instrukčních sad jsou:
 \begin{itemize}
  \item \alert{Řídící}
  \begin{itemize}
   \item \alert{BRANCH} (též \alert{JUMP}) -- přesuň se na jinou adresu v paměti
    a vykonej instrukce tam;
   \pause
   \item \alert{CONDITIONALLY BRANCH} -- je-li daná podmínka splněna, přesuň se
    na jinou adresu v paměti a vykonej instrukce tam;
   \pause
   \item \alert{INDIRECTLY BRANCH} -- přečti adresu z daného místa
    (paměť, registr, \dots) a vykonej instrukce na této adrese;
   \pause
   \item \alert{CALL} -- proveď instrukce mezi danými adresami v paměti, po
    vykonání se vrať zpět a pokračuj další instrukcí.
  \end{itemize}
 \end{itemize}
\end{frame}

\section{Programovací jazyky}

\begin{frame}
 \frametitle{Co je programovací jazyk?}
 Programovací jazyk je způsob zápisu počítačového programu, který je blíže
 lidskému jazyku než program psaný čistě v CPU instrukcích.
 \pause
 Programovací jazyk (stejně jako lidský) je popsán svoji
 \begin{itemize}
  \item \alert{formou/syntaxí} -- způsobem zápisu,
  \pause
  \item \alert{sémantikou} -- významem slov (posloupností symbolů).
 \end{itemize}
 \pause
 Každý programovací jazyk potřebuje buď svůj (aspoň jeden) \alert{překladač}
 (compiler) nebo \alert{interpret} (interpreter), aby v něm psané programy mohly
 být vykonány CPU jako sady základních instrukcí.
\end{frame}

\begin{frame}[fragile]
 \frametitle{Syntax programovacího jazyka}
 Syntax programovacího jazyka je obvykle definována přes \alert{regulární
 výrazy} (pro lexikální strukturu) a tzv. \alert{Backusovou-Naurovou formou}
 (pro gramatickou strukturu).
 \pause
 Příklad syntaxe velmi jednoduchého programovacího jazyka může vypadat takto:
 \begin{verbatim}
  expression ::= atom | list\\
  atom       ::= number | symbol\\
  number     ::= [+-]?['0'-'9']+\\
  symbol     ::= ['A'-'Z''a'-'z'].*\\
  list       ::= '(' expression* ')'
 \end{verbatim}
\end{frame}

\begin{frame}[fragile]
 \frametitle{Sémantika programovacího jazyka}
 Významovou stránku programovacího jazyka je obecně mnohem těžší formalizovat.
 \pause
 Obvykle se skládá ze dvou částí:
 \begin{itemize}
  \item \alert{statická sémantika} -- omezení na strukturu \uv{vět}, která je
   obtížné nebo nemožné definovat v rámci syntaxe.\\
   \pause
   U většiny programovacích jazyků například udává, že proměnné musejí být
   inicializovány před použitím nebo že dvě různé větve stejné podmínky musejí
   být podmíněny různými výroky.
 \end{itemize}
\end{frame}

\begin{frame}[fragile]
 \frametitle{Sémantika programovacího jazyka}
 Významovou stránku programovacího jazyka je obecně mnohem těžší formalizovat.
 Obvykle se skládá ze dvou částí:
 \begin{itemize}
  \item \alert{dynamická sémantika} -- definuje způsob, jakým se mají
   interpretovat (v CPU instrukcích) konkrétní věty.\\
   \pause
   Například udává, že
   \begin{verbatim}
    if x < 0:
        print("x is less than 0.")
   \end{verbatim}
   \vspace*{-1.5em}
   má podmínečně skočit, kde je uložen program \texttt{print("x is less than
   0.")}, vykonat ho, pak se vrátit a pokračovat další instrukcí.
 \end{itemize}
\end{frame}

\begin{frame}
 \frametitle{Datové typy}
 Každé datum, které ukládáte přes programovací jazyk, má konkrétní \alert{typ}
 -- např. celé číslo, desetinné číslo, string, pole apod.\\
 \pause
 Datové typy jsou pro moderní programovací jazyky v podstatě nezbytné, protože
 zjednodušené syntaxe je dosaženo na úkor přesnosti vyjádření instrukcí.\\
 \pause
 Mnoho programovacích jazyků již docílilo kompletní automatizace správy paměti
 (např. JavaScript, Python, C\# apod.)
\end{frame}

\begin{frame}
 \frametitle{Datové typy}
 Programovací jazyky můžeme rozdělit do kategorií podle toho, jak zacházejí s
 datovými typy:
 \pause
 \begin{itemize}
  \item \alert{netypované jazyky} -- programovací jazyky \alert{bez datových
  typů}, programátor tedy provádí operace přímo na uložená data (posloupnosti 1
  a 0). Sám musí vědět, co daná posloupnost znamená. V podstatě pouze nejstarší
  jazyky, nebo obecně jazyky velmi blízko strojovému kódu (instrukce psané přímo
  v 0 a 1), jsou netypované.
 \pause
 \item \alert{typované jazyky} -- programovací jazyky, kde každé v paměti
  uložené datum má jasně přiřazený typ. Datový typ bývá obvykle uložen v
  několika prvních bitech samotného data. To znamená, že třeba posloupnost
  \texttt{01001011} bude v programovacím jazyku interpretována podle toho, jaké
  bity jí předchází.
 \end{itemize}
\end{frame}

\begin{frame}
 \frametitle{Typované jazyky}
 Typované programovací jazyky pak můžeme ještě rozdělit podle toho, jestli nutí
 programátora typy deklarovat či nikoliv.
 \pause
 \begin{itemize}
  \item \alert{staticky typované} jazyky vyžadují u inicializace každé proměnné
   explicitně vypsat její typ. Jejich syntaxe zakazuje používat operace na
   datové typy, pro které nebyly definovány.
  \pause
  \item \alert{dynamicky typované} jazyky nevynucují explicitní vypsání typu a
   často se typ proměnné za běhu programu mění. Kompatibilita operací s typem
   proměnné se ověřuje až za překladu nebo interpretace.
 \end{itemize}
\end{frame}

\begin{frame}
 \frametitle{Rozšířené programovací jazyky}
 \begin{enumerate}
  \item \alert{Python} (dynamicky typovaný, překládaný i interpretovaný) --
   víceúčelový programovací jazyk. Používá se hlavně na
   \begin{itemize}
    \item datovou analýzu a vývoj umělé inteligence,
    \item vývoj webových aplikací,
    \item testování a simulaci chování softwaru.
   \end{itemize}
  \pause
  \item \alert{Java} (staticky typovaný, překládaný i interpretovaný) --
   používaný hlavně na
  \begin{itemize}
   \item vývoj mobilních aplikací,
   \item vývoj cloudových služeb,
   \item chatboty a marketing,
   \item rozsáhlé webové aplikace.
  \end{itemize}
 \end{enumerate}
\end{frame}

\begin{frame}
 \frametitle{Rozšířené programovací jazyky}
 \begin{enumerate}
  \setcounter{enumi}{2}
  \item \alert{JavaScript} (dynamicky typovaný, interpretovaný) -- hlavní
   využití JS je ve vývoji webových aplikací (konkrétně frontendu = toho, co vidí
   uživatel).
  \pause
  \item \alert{C/C++} (staticky typovaný, překládaný) -- extrémně rychle
   překládaný jazyk, používaný hlavně na
  \begin{itemize}
   \item vývoj překladačů jiných jazyků (překladá například Python),
   \item vývoj operačních systémů,
   \item vývoj databází,
   \item vývoj aplikací.
  \end{itemize}
 \end{enumerate}
\end{frame}

\end{document}
